\documentclass[11pt]{article}
\usepackage[margin=1in]{geometry}
\usepackage{amsmath,amssymb,booktabs}
\title{Degree-five and degree-six\\quintic mirror coefficients}
\author{The Last Math Competition, conjecture 00000001821}
\date{}
\begin{document}
\maketitle

The quintic B-model predicts these instanton numbers in degrees five and six
\[
n_5=229{,}305{,}888{,}887{,}625,\qquad n_6=248{,}249{,}742{,}118{,}022{,}000.
\]
The values in the conjecture belong to degrees three and four. We compute
the mirror coefficients directly, which disproves the asserted numerical
agreement with the B-model.

\section*{Formal coefficient formula}

Use the large-complex-structure coordinate \(z=(5\psi)^{-5}\). Put
\[
 A_j=\frac{(5j)!}{(j!)^5},\qquad
 I_0(z)=\sum_{j\ge 0} A_jz^j,\qquad
 I_1(z)=5\sum_{j\ge 1} A_j\left(\sum_{m=j+1}^{5j}\frac1m\right)z^j.
\]
Here \(I_1\) is the regular part of the logarithmic period. Set
\(G=I_1/I_0\), so the flat coordinate is \(q=z\exp G(z)\).
In this coordinate the normalized B-model Yukawa coupling is
\begin{equation}
 K(q(z))=\frac{5}{(1-3125z)I_0(z)^2(1+zG'(z))^3}.
 \label{eq:yukawa}
\end{equation}
This is the quintic period and Yukawa formula of Candelas, de la Ossa,
Green, and Parkes \cite[sections~4.2 and~5.3]{cogp}, with constant term
\(K(0)=5\). The normalization by \(I_0^{-2}\) and the cube of the
logarithmic coordinate derivative give the last two factors of the denominator.

For \(d\ge 1\), write \(c_d=[q^d]K(q)\). Formal residue substitution and
\(q'(z)/q(z)=z^{-1}(1+zG'(z))\) give
\begin{equation}
 c_d=[z^d] \frac{5\exp(-dG(z))}
                         {(1-3125z)I_0(z)^2(1+zG'(z))^2}.
 \label{eq:residue}
\end{equation}
Indeed, substitute \(q=q(z)\) into \(K(q)q^{-d-1}\,dq\) and take the
coefficient of \(z^{-1}\,dz\). Residue is invariant in this substitution
because for each monomial differential \(q^m\,dq\), \(m\ne-1\), the
substituted differential is the derivative of \(q(z)^{m+1}/(m+1)\), with
zero residue. The remaining case \(dq/q\) has residue one. Since
\(q(z)=z+O(z^2)\), only finitely many terms contribute to the residue.

\section*{Exact extraction}

The instanton numbers \(n_d\) are defined by the Lambert series
\[
 K(q)=5+\sum_{d\ge1}d^3n_d\frac{q^d}{1-q^d}.
\]
Expanding each geometric series gives \(c_d=\sum_{e\mid d}e^3 n_e\).
Thus \(d^3n_d=c_d-\sum_{e\mid d,\ e<d}e^3n_e\), a triangular recurrence.
To compute equation~\eqref{eq:residue} to degree six, truncate every series
at degree six. For a unit series \(a\), its reciprocal \(b\) satisfies
\(b_0=1/a_0\) and \(a_0b_j=-\sum_{i=1}^j a_i b_{j-i}\) for \(j\ge1\).
For \(E=\exp(B)\) with \(B(0)=0\), the identity \(E'=B'E\) gives
\(E_0=1\) and \(jE_j=\sum_{i=1}^j iB_iE_{j-i}\). Every coefficient on
the right is of degree at most the requested degree, so the truncation is
exact for the entries in the table.

\begin{center}
\begin{tabular}{rrr}
\toprule
Degree \(d\) & Yukawa coefficient \(c_d\) & Instanton number \(n_d\) \\
\midrule
1 & \(2{,}875\) & \(2{,}875\) \\
2 & \(4{,}876{,}875\) & \(609{,}250\) \\
3 & \(8{,}564{,}575{,}000\) & \(317{,}206{,}375\) \\
4 & \(15{,}517{,}926{,}796{,}875\) & \(242{,}467{,}530{,}000\) \\
5 & \(28{,}663{,}236{,}110{,}956{,}000\) & \(229{,}305{,}888{,}887{,}625\) \\
6 & \(53{,}621{,}944{,}306{,}062{,}201{,}000\) & \(248{,}249{,}742{,}118{,}022{,}000\) \\
\bottomrule
\end{tabular}
\end{center}

The instanton numbers agree with Table~4 of the original mirror
calculation \cite{cogp}. The Lean project checks both columns by exact
rational arithmetic from the periods, the residue formula, and the divisor
recurrence. Its two inequality theorems show that both degree assignments
in the conjecture fail. This calculation addresses the B-model Lambert
coefficients independently of the geometric interpretation of each
number as a count of rational curves.

\begin{thebibliography}{1}
\bibitem{cogp}
P.~Candelas, X.~C. de la Ossa, P.~S. Green, and L.~Parkes,
\emph{A pair of Calabi--Yau manifolds as an exactly soluble superconformal
theory}, Nuclear Physics B \textbf{359} (1991), 21--74.
\texttt{doi:10.1016/0550-3213(91)90292-6}.
\end{thebibliography}
\end{document}
